\chapter{Scope, historical evidence and method}\label{ch:scope}
\plain{The equations, the stored history and the picture answer different
questions. A picture can show a line without telling us who arrived first.
A program can define a clock without counting every time an animation callback
changes the state. We reconstruct these distinctions before interpreting the motion.}

\section{The corpus and the scientific question}
This volume asks what three recovered moving-center simulators actually
calculated and displayed, and what their surviving evidence establishes.
The primary programs are designated S01, S02 and S03 in research order
(Table~\ref{tab:corpus}). These identifiers do not assert a historical
creation sequence. The accepted construction studies and the separately
authorized three-point replay supply the basis for this synthesis
\cite{construction01,replay01,construction02,study03,charter}.
\claim{C01}

\begin{table}[htbp]\centering\small
\caption{The selected corpus. Program identity is fixed by the SHA-256 records
in Appendix~\ref{app:provenance}; media association does not itself certify
an exact generating run.}\label{tab:corpus}
\begin{tabularx}{\linewidth}{@{}l Y Y@{}}\toprule
ID and source & Native construction & Surviving evidence used\\\midrule
S01: recursive\_attractor.py & Three 2D points, pairwise response,
zero initial velocities, 500 sweeps & Attributed historical static PNG;
separate seed-0 reproduction with full saved chronology\\
S02: 3Dtriangles.py & Eight 3D centers, sequential centroid response,
random initial velocities & Attributed 2,000-frame historical movie;
no certified numerical history for that movie\\
S03: plot\_1.py & Same centroid family plus eight binary response modes,
stochastic toggles and diagnostics & Attributed 200-frame movie; original
logs A and B; attributed historical body-7 charts\\\bottomrule
\end{tabularx}
\end{table}

The scope is deliberately narrower than a general SRG theory. We do not merge
these programs with fixed-grid systems, later amplitude/neighbor-change
models, production physics or memory kernels. Names such as REFU, TGMO, RPCO
and GRF are read through their actual implementations. A variable named
\filepath{force} is not by itself a physical-force law; a function labelled
“memory” need not read stored routes. No Newtonian or relativistic field
equation, new metric, additional potential or hidden following rule is
assumed. \claim{C02}

\section{Evidence classes and attribution}
We use the charter's distinctions throughout:
\begin{description}[style=nextline,leftmargin=0pt]
\item[ORIGINAL] Preserved historical bytes: programs, figures, movies, logs
or contemporaneous records.
\item[ATTRIBUTED] A relationship between an artifact and a source or run
supported by corroboration, with the remaining uncertainty stated.
\item[OBSERVED] A directly inspected or quantified feature of specified
original or reproduced evidence.
\item[REPRODUCED] A new execution of unchanged rules, identified by its
initializer and environment; it is not relabelled as lost original data.
\item[DERIVED] A mathematical conclusion from explicit source equations and
stated conditions, or an analysis product constructed from existing records.
\item[HYPOTHESIS] An interpretation or research lead not established by the
preceding evidence.
\item[MISSING / AMBIGUOUS] An unrecovered identity, state, dependency,
chronology or mechanism. Missing evidence is not silently supplied.
\end{description}

These labels can coexist. A movie is an original artifact even when its link
to a particular source snapshot is only attributed. Matching a filename,
title, geometry and nominal frame rate strengthens source association;
it does not recover a random seed, callback sequence or export receipt.
Likewise, a derived figure may display original measurements without becoming
an original historical plot. \claim{C03}

No original script sets a fixed seed. S01's source makes no file export at
all; its historical PNG therefore has an unrecovered saving step. S02 and
S03 write fixed movie names relative to the process working directory.
Those names can be reused by later executions. S03 gives its CSV a
second-resolution timestamp, without a timezone or unique run identifier.
The \filepath{PLOT_1_} prefix on log A is not produced by that filename
expression. Neither current directory placement nor filename dates alone
establish which source version generated an artifact.
See \src{S01}{1--74}, \src{S02}{110--119} and \src{S03}{74--78,158--167}.
\claim{C04}

\section{Method and limits of the present synthesis}
For the underlying studies, the complete programs were read, their state and
execution order were transcribed, relevant original evidence was inspected,
and numerical consistency was checked at the declared coverage. This volume
reuses those accepted results and their saved verification receipts.
It does not rerun the original simulations or their scientific checkers.
New operations are editorial: preserving source copies, checking hashes and
links, assembling figures, typesetting, inspecting the PDF and packaging.
\claim{C05}

For S01's accepted replay, all 500 recorded end-of-sweep position arrays and
the final velocity array matched an execution of the unchanged copied source
bit for bit in the recorded environment. That comparison does not certify
the lost historical PNG or promise bitwise identity on every software stack.
For S03, all-row energy, position and incoming-GRF checks coexist with
\emph{sampled} full velocity-equation checks. The latter must not be described
as an exhaustive audit of every force calculation.
\cite{replay01,study03} \claim{C06}

The canonical study ledger is a navigation record. Its original SIM03 entry
was captured while that study was in review. The Volume I work order records
GPT's subsequent acceptance within the evidence/provenance qualifications.
The ledger is minimally updated accordingly; the accepted study package and
its historical in-review snapshot remain unchanged. Scientific and editorial
review is complete; publication requires separate authorization. \claim{C07}

\section{Notation, units and record locators}
Positions and velocities are numerical arrays. The programs supply no
physical calibration of coordinate length, mass or time. We use
“coordinate units” and “model-step units” rather than metres or seconds
for the dynamics. Video seconds are actual encoded presentation times;
they are not silently identified with model time.

For S01, $n=0,\ldots,499$ is the source loop index and $S_n$ the state after
$n$ completed sweeps. Event $e=3n+i+1$ updates body $i$ in sweep $n+1$.
For S03, $q$ is a zero-based callback-block index introduced by the accepted
analysis in physical CSV order; $f_q$ is the logged \filepath{Step} label.
Body $i$ at callback $q$ is on physical CSV line $8q+2+i$, counting the header
as line~1. This locates repeated labels without deduplication. Movie frame
indices are zero-based throughout. \claim{C08}

Source locators use the unchanged copies, whose full originals and SHA-256
identities are listed in Appendix~\ref{app:provenance}. Major claims carry
matrix IDs such as \claim{C08}; the accompanying
\filepath{CLAIM_EVIDENCE_MATRIX.md} maps each to exact support and caveats.
Figures have independent provenance in
\filepath{FIGURE_AND_MEDIA_REGISTER.md}. These registers are part of the
scientific review package, not a substitute for the explanations below.
